\documentclass[12pt,letterpaper]{report}
\usepackage{graphicx}
\usepackage{scrextend}
\usepackage{vmargin}
\usepackage{graphicx}
\usepackage{multirow}
\usepackage[utf8]{inputenc}
\usepackage[spanish]{babel}
\usepackage{multicol}
\usepackage{enumerate}
\usepackage{float}
\usepackage{amsmath, amsthm, amssymb, amsfonts}
\usepackage[usenames]{color}
\parindent=0mm
\pagestyle{empty}
\definecolor{miorange}{rgb}{0.91, 0.43, 0.0}
\begin{document}
\setmargins{2.5cm}      
{1.5cm}                     
{2cm}  
{24cm}                    
{10pt}                          
{1cm}                          
{0pt}                             
{2cm}
\begin{flushright}
\today
\end{flushright}
\vspace{-1.75cm}
\section*{Tarea Clase 6}
\subsection*{Obtenga las soluciones de las siguientes ecuaciones:}
\begin{enumerate}
\item \begin{equation*}
    2x-(5x+3)=7+(3x-2)
\end{equation*}
\begin{align*}
    2x-5x-3&=7+3x-2\\
    -3x-3&=5+3x\\
    6x&=-8\\
    x&=-\frac{8}{6}
\end{align*}
\item \begin{equation*}
    5x-9=3(x-2)
\end{equation*}
\begin{align*}
    5x-9&=3x-6\\
    2x&=15\\
    x&=15/2
\end{align*}
\item \begin{equation*}
    -2+x-8=2-x+8
\end{equation*}
\begin{align*}
    2x&=20\\
    x&=10
\end{align*}
\item \begin{equation*}
    2x+x+5=18-3x
\end{equation*}
\begin{align*}
    6x&=13\\
    x&=13/6
\end{align*}
\item \begin{equation*}
    -2x-5=8x+3
\end{equation*}
\begin{align*}
    10x&=-8\\
    x&=-8/10
\end{align*}
\item \begin{equation*}
    2x+3(x-1)=6-4(2x+3)
\end{equation*}
\begin{align*}
    2x+3x-3&=6-8x-12\\
    13x&=-3\\
    x&=-3/13
\end{align*}
\item \begin{equation*}
    3x+5=\frac{5}{2}x+4
\end{equation*}
\begin{align*}
    \frac{1}{2}x&=-1\\
    x&=-2
\end{align*}
\item \begin{equation*}
    4(y+1)+3y-1=7y+3
\end{equation*}
\begin{align*}
    4y+4+3y-1&=7y+3\\
    7y+3&=7y+3\\
    0&=0
\end{align*}
este problema tiene soluciones infinitas
\item \begin{equation*}
    8x-15x-30x-51x=53x+31x-172
\end{equation*}
\begin{align*}
    -172x&=-172\\
    x&=1
\end{align*}
\item \begin{equation*}
    \frac{3}{2}x+\frac{1}{4}+2=\frac{3}{4}x-\frac{1}{3}x
\end{equation*}
\begin{align*}
    \frac{13}{12}x&=-\frac{5}{2}
    x&= -\frac{30}{13}
\end{align*}
\item \begin{equation*}
    x^2+6=7x
\end{equation*}
\begin{align*}
    x^2-7x+6&=0\\
    (x-6)(x-1)&=0\\
    x&=1,6
\end{align*}
\item \begin{equation*}
    3x^2+2x-8=0
\end{equation*}
\begin{align*}
    (3x-2)(x+4)&=0\\
    x&=-4,2/3
\end{align*}
\item \begin{equation*}
    2x^2-12x+18=0
\end{equation*}
\begin{align*}
    x^2-6x+9&=0\\
    (x-3)^2&=0\\
    x&=3,3
\end{align*}
\item \begin{equation*}
    -x^2=-3x-10
\end{equation*}
\begin{align*}
    x^2-3x-10&=0\\
    (x-5)(x+2)&=0\\
    x&=5,-2
\end{align*}
\item \begin{equation*}
    x^2-2x=1
\end{equation*}
\begin{align*}
    x^2-2x-1&=0\\
\end{align*}
este se resuelve con formula general donde, a=1,b=-2,c=-1
\item \begin{equation*}
    -3x^2-6x+9=0
\end{equation*}
\begin{align*}
    x^2+2x-3&=0\\
    (x+3)(x-1)&=0\\
    x&=1,-3
\end{align*}
\item \begin{equation*}
    -2x^2-10x=0
\end{equation*}
\begin{align*}
    x(x+5)&=0\\
    x&=0,-5
\end{align*}
\item \begin{equation*}
    3x^2=-6x-3
\end{equation*}
\begin{align*}
    3x^2+6x+3&=0\\
    x^2+2x+1&=0\\
    (x+1)^2=0\\
    x=-1,-1
\end{align*}
\end{enumerate}
\end{document}